% 第三批：《大学物理简明教程》第5版，解答题3.3.1--3.3.8。
% 源书/PDF第70--72页。公式独立复算，原书图题模型单列，待二校。
\subsection{解答题：3.3.1--3.3.8}

\sourceproblem{3.3.1}{70}{70}
\problemtext 半径$R=0.30\,\mathrm m$的飞轮由静止开始做匀角加速度$\alpha=0.5\,\mathrm{rad/s^2}$的转动。求轮缘一点转过$240^\circ$时的角加速度、角速度及切向加速度。
\workings
$\alpha$恒定，故角加速度仍为$0.5\,\mathrm{rad/s^2}$。
$\theta=240^\circ=4\pi/3$，运动学关系$\omega^2=2\alpha\theta$给
\[
\boxed{\omega=\sqrt{\frac{4\pi}{3}}\approx2.047\,\mathrm{rad/s}},
\qquad \boxed{a_\tau=R\alpha=0.15\,\mathrm{m/s^2}}.
\]
\crosscheck 若另外需要法向加速度，则$a_n=R\omega^2=0.4\pi\,\mathrm{m/s^2}$；题面只要求角加速度、角速度及切向加速度。

\sourceproblem{3.3.2}{70--71}{70--71}
\problemtext 半径$R=10\,\mathrm{cm}$、内壁光滑的半球形碗绕竖直对称轴匀速转动。一个小球相对碗静止，位置高于碗底$h=4\,\mathrm{cm}$。求所需角速度。
\workings
设小球到轴的水平距离为$r$、到球心的竖直投影绝对值$R-h$。
碗壁支持力$N$沿指向球心的法线，竖直方向平衡、水平向心：
\[
N\frac{R-h}{R}=mg,\qquad N\frac rR=m\omega^2r.
\]
排除$r=0$的底端特例，两式相除得
\[
\boxed{\omega=\sqrt{\frac{g}{R-h}}}
=\sqrt{\frac{9.8}{0.10-0.04}}
\approx12.78\,\mathrm{rad/s}.
\]
\crosscheck 位置几何$r^2=2Rh-h^2=(0.08\,\mathrm m)^2>0$；
当$h\to R^-$时所需角速度发散。

\sourceproblem{3.3.3}{70--71}{70--71}
\problemtext 均匀圆盘质量$m=1.1\,\mathrm{kg}$、转动惯量$J=\frac12mr^2$，
轮缘绳悬挂$m_1=1.0\,\mathrm{kg}$重物。施加外力矩使重物以$v_0=0.6\,\mathrm{m/s}$匀速上升，突然撤去外力矩，求重物开始反向运动前所经时间。
\workings
撤去外力矩后，重力使圆盘角速度减小。以重物向上为正、受力$T-m_1g=m_1a$，轮轴上由绳张力引起的减速力矩满足$-Tr=J(a/r)$。
联立或由整体有效质量得到
\[
a=-\frac{m_1g}{m_1+J/r^2}
=-\frac{m_1g}{m_1+m/2}.
\]
至速度为零所需时间
\[
\boxed{t=\frac{v_0}{|a|}
=\frac{v_0(m_1+m/2)}{m_1g}
=\frac{0.6(1+1.1/2)}{9.8}\approx0.0949\,\mathrm s}.
\]
\crosscheck $r$从有效转动惯量中消去，因圆盘惯量正比$r^2$；撤矩后的初速度仍向上，随后由重力制动。

\sourceproblem{3.3.4}{71}{71}
\problemtext 半径$R$、转动惯量$J$的水平光滑转台，以角速度$\omega_0$转动。质量$m$的人最初站在轴心，随后沿半径走到边缘，求新的角速度。
\workings
对竖直轴外力矩忽略，角动量守恒。
起初人的质点转动惯量可忽略，后为$mR^2$，因此
\[
J\omega_0=(J+mR^2)\omega_f,
\qquad \boxed{\omega_f=\frac{J}{J+mR^2}\omega_0}.
\]
\crosscheck $m\to0$恢复$\omega_0$；人移到更远处使系统角速率降低；非保守内功使动能未必守恒。

\sourceproblem{3.3.5}{71}{71}
\problemtext 质量$m=60\,\mathrm{kg}$、半径$R=0.25\,\mathrm m$的均匀圆盘飞轮以$900\,\mathrm{r/min}$旋转。杠杆左端为支点，支点到闸瓦接触点为$0.50\,\mathrm m$，闸瓦到施力端为$0.75\,\mathrm m$；闸瓦与轮缘动摩擦系数$\mu=0.4$。$(1)$若向下施力$F=100\,\mathrm N$，求停转时间与圈数；$(2)$若要求$2\,\mathrm s$内转速减半，求$F$。
\workings
飞轮转动惯量$J=\frac12mR^2=1.875\,\mathrm{kg\,m^2}$，
初角速度$\omega_0=900(2\pi)/60=30\pi\,\mathrm{rad/s}$。
闸杆视为无质量杠杆，关于左端支点力矩平衡
\[
N(0.50)=F(0.50+0.75),\qquad N=2.5F.
\]
摩擦制动力矩大小
$\tau=\mu NR=\mu(2.5F)R=0.25F\,\mathrm{N\,m}$
（$F$以牛顿的数值计）。
$(1)$ $F=100\,\mathrm N$时，
\[
|\alpha|=\tau/J=25/1.875=\frac{40}{3}\,\mathrm{rad/s^2},
\quad
\boxed{t_{\rm stop}=\frac{\omega_0}{|\alpha|}
\approx7.069\,\mathrm s}.
\]
停转角位移$\theta=\omega_0^2/(2|\alpha|)$，故
\[
\boxed{n=\frac{\theta}{2\pi}
=\frac{\omega_0^2}{4\pi|\alpha|}
\approx53.01\text{圈}}.
\]
$(2)$ $2$秒降到$\omega_0/2$要求
$|\alpha|=\omega_0/(4\,\mathrm s)$，于是
\[
\boxed{F=\frac{J\omega_0}{4(0.25)(2\,\mathrm s)}
=\frac{1.875(30\pi)}{1.0(2)}
\approx176.7\,\mathrm N}.
\]
\crosscheck 第二问的简化公式直接写为
$F=J[\omega_0/(2\Delta t)]/(\mu\,2.5R)$；
上式数值按$\Delta t=2\,\mathrm s$代入，见此式。
两个力都大于零，制动时间与制动力成反比。
\sourceproblem{3.3.6}{71}{71}
\problemtext 两个同轴均匀圆柱体固连，可绕水平轴无摩擦转动。大、小圆柱分别质量$M=10\,\mathrm{kg},m=4\,\mathrm{kg}$，半径$R=0.20\,\mathrm m,r=0.10\,\mathrm m$；绳分别在两柱绕向相反方向，吊质量$m_1=m_2=2\,\mathrm{kg}$的重物。求角加速度和两绳张力。
\workings
取$m_1$向下为正，$m_2$向上为正。合转动惯量
\[
J=\tfrac12MR^2+\tfrac12mr^2=0.22\,\mathrm{kg\,m^2}.
\]
加速度大小$a_1=\alpha R,\ a_2=\alpha r$，牛顿方程
\[
m_1g-T_1=m_1R\alpha,\quad
T_2-m_2g=m_2r\alpha,\quad
T_1R-T_2r=J\alpha.
\]
消去张力得
\[
\boxed{\alpha=\frac{m_1gR-m_2gr}
{J+m_1R^2+m_2r^2}
=\frac{2(9.8)(0.2-0.1)}{0.22+0.08+0.02}
=6.125\,\mathrm{rad/s^2}}.
\]
从而
\[
\boxed{T_1=m_1(g-R\alpha)=17.15\,\mathrm N},\qquad
\boxed{T_2=m_2(g+r\alpha)=20.825\,\mathrm N}.
\]
\crosscheck 直接代入$T_1R-T_2r=0.22\alpha$两边均为$1.3475\,\mathrm{N\,m}$。

\sourceproblem{3.3.7}{71--72}{71--72}
\problemtext 两个半径均为$R$、质量均为$m$的均匀圆盘状定滑轮，轻绳绕过两轮，两端悬挂质量分别为$m$和$2m$的重物。绳轮不滑动，求两滑轮间绳段的张力。
\workings
取$2m$侧向下为正。两轮转动惯量各为$J=\frac12mR^2$，整套系统有效惯量质量
$m+2m+2(J/R^2)=4m$。
\[
a=\frac{(2m-m)g}{4m}=\frac g4.
\]
设左侧张力$T_L=m(g+a)=\frac54mg$，右侧张力
$T_R=2m(g-a)=\frac32mg$。
左轮由中间张力$T_C$与左侧张力驱动：
$(T_C-T_L)R=Ja/R$，因此
\[
\boxed{T_C=T_L+\frac{Ja}{R^2}
=\frac54mg+\frac18mg=\frac{11}{8}mg}.
\]
\crosscheck 由右轮$(T_R-T_C)R=Ja/R$也得
$T_C=\frac32mg-\frac18mg=\frac{11}{8}mg$。

\sourceproblem{3.3.8}{71--72}{71--72}
\problemtext 长度$l$、质量$m$的均匀细杆在竖直平面内可绕中点$O$自由转动，初始水平且静止。一只质量同为$m$的小虫以垂直于细杆的速率$v_0$落在距$O$为$l/4$处，并沿杆向端点$A$爬行。求使杆角速度保持常值时，小虫需要的沿杆爬行速率。
\workings
\textbf{必须区分“恒定角速度”和“恒定爬行速率”，原题只要求前者。}
初始碰撞过程中对$O$轴角动量守恒：
\[
m v_0\frac l4
=\left(\frac{ml^2}{12}+m\frac{l^2}{16}\right)\omega_*,
\qquad \boxed{\omega_*=\frac{12v_0}{7l}}.
\]
碰后取由水平向下转过的角度$\theta=\omega_* t$为正，小虫距轴$r(t)$沿正方向向外，杆的质心位于轴上而重力矩为零。对整体应用角动量定理：
\[
\frac{d}{dt}\left[\left(\frac{ml^2}{12}+mr^2\right)\omega_*\right]
=mg r\cos\theta.
\]
因$\dot\omega_*=0$，只要$r>0$，可约去$mr$，得
\[
\boxed{v_{\rm crawl}(t)=\dot r
=\frac{g}{2\omega_*}\cos(\omega_*t)}.
\]
这给出\emph{保持角速度恒定所需的变速爬行规律}。
积分得
$r(t)=l/4+\frac{g}{2\omega_*^2}\sin(\omega_*t)$。
\crosscheck 在$0\le\theta\le\pi/2$时$\dot r\ge0$，可向$A$端爬行；
越过$\pi/2$后要求$\dot r<0$，故不能无限时间保持“持续向外爬行”和“恒定角速度”两项条件同时成立。若题目另要求爬到端点，还需检验$r(t)\le l/2$和初速度对应的可达条件。
